\maketitle

\begin{abstract}
A fitted optical intercept cannot be recovered uniquely from a published
HgCdTe absorption curve until the extraction equation, fit-domain
endpoints, and point-weighting convention are declared. This
two-spectrum case study quantifies that downstream dependence for 300 K
absorption-coefficient curves reported by Moazzami \emph{et al.} at
nominal compositions \(x=0.226\) and \(x=0.310\). The source curves were
themselves obtained by infrared spectroscopic ellipsometry and
optical-model inversion; the present records were reconstructed from
embedded bitmap figures rather than native instrument exports. The
analysis is therefore conditional on the published optical inversion and
does not estimate the total experimental uncertainty of either specimen.

The dominant immediately observable sensitivity is the fit convention.
Across a predeclared grid of lower and upper absorption endpoints and
three point-weighting rules, the free-exponent fitted intercept moves by
as much as 5.71 meV for \(x=0.226\) and 5.44 meV for \(x=0.310\).
Alternative deterministic reconstructions with nominal fit membership
held fixed move fitted intercepts by at most 0.69 and 0.81 meV; allowing
reconstruction-dependent membership increases the largest admissible
shifts to 1.25 and 1.96 meV. Coherent calibration-corner perturbations
produce maxima of 0.89 and 0.68 meV. A visible two-point reconstruction
irregularity near 0.198 eV changes a non-boundary fitted intercept by
approximately 0.12 meV when omitted and is not the controlling feature.

The primary model comparison is restricted to fitted-intercept forms:
free and fixed fractional-power relations, source-panel exponents, and a
provenance-bounded Chu Kane-region relation. Fixed-absorption crossings
are retained separately as operational coordinates. After
boundary-limited fits are excluded, the complete declared fitted-model
spans are 5.09 and 6.68 meV. Energy-resolved residuals show systematic
model-dependent structure, and several fixed forms do not reproduce the
reconstructed curve within its apparent line-width resolution. A
residual-envelope screen is reported only as a secondary descriptive
diagnostic; its threshold sensitivity is shown explicitly and it is not
used as a confidence interval or primary headline.

The Hansen-Seiler difference is used only as a numerical scale, not as
an equation ranking, because specimen-level composition uncertainty is
unavailable and the printed composition precision alone is insufficient
for a sub-meV comparison. The defensible conclusion is conditional: for
these two reconstructed, model-derived spectra, the fitted optical
intercept is materially conditioned by extraction model and especially
by fit-domain and weighting choices. The analysis does not identify a
unique physical gap, a universal correction, or a preferred universal
HgCdTe gap equation.
\end{abstract}

\noindent\textbf{Keywords:} HgCdTe; optical absorption; fitted optical intercept; model-form dependence; fit-window sensitivity; spectroscopic ellipsometry; reconstruction sensitivity

\vspace{0.8em}
